\section{Merancang Grammar Subset-C}
\label{sec:ch08-subsetc}

\subsection{Prinsip Perancangan}

Tata bahasa Subset-C dirancang dengan prinsip berikut:
\begin{enumerate}[label=\arabic*.,leftmargin=*]
  \item satu nonterminal untuk satu tingkat prioritas ekspresi;
  \item operator dengan prioritas lebih tinggi berada lebih dalam
        pohon;
  \item pernyataan dan deklarasi dipisahkan;
  \item bagian opsional ditandai dengan nonterminal \code{opt\_*};
  \item pengulangan ditandai dengan nonterminal \code{*\_tail}.
\end{enumerate}

\subsection{Struktur Nonterminal}

\begin{figure}[H]
\centering
\begin{tikzpicture}[node distance=5mm and 6mm]
  \node[phase, text width=2.4cm] (program) {program};
  \node[phase, text width=2.4cm, below left=of program] (top) {topdecl};
  \node[phase, text width=2.4cm, below=of program] (stmt) {stmt};
  \node[phase, text width=2.4cm, below right=of program] (expr) {expr};
  \node[phase, text width=2.4cm, below=of expr] (assign) {assign};
  \node[phase, text width=2.4cm, below=of assign] (logor) {logor};
  \draw[arrow] (program) -- (top);
  \draw[arrow] (program) -- (stmt);
  \draw[arrow] (program) -- (expr);
  \draw[arrow] (expr) -- (assign);
  \draw[arrow] (assign) -- (logor);
\end{tikzpicture}
\caption{Hierarki nonterminal utama Subset-C}
\label{fig:ch08-hierarki}
\end{figure}

\subsection{Deklarasi dan Fungsi}

Aturan untuk deklarasi fungsi:
\[
\begin{aligned}
\text{topdecl} &\to \text{type}\ \code{ID}\ \text{"("}\
\text{opt\_params}\ \text{")"}\ \text{topdecl\_tail},\\
\text{topdecl\_tail} &\to \text{block} \mid \text{";"}.
\end{aligned}
\]
Aturan ini memungkinkan deklarasi fungsi dengan blok tubuh atau
deklarasi prototipe tanpa tubuh. Parameter opsional ditangani oleh
\code{opt\_params}.

\subsection{Pernyataan}

Aturan pernyataan mengelompokkan konstruksi kontrol:
\[
\begin{aligned}
\text{stmt} &\to \text{decl}\ \text{";"} \mid \text{block}
\mid \text{ifstmt} \mid \text{whilestmt}\\
&\quad \mid \text{forstmt} \mid \text{returnstmt}
\mid \text{breakstmt} \mid \text{continuestmt}\\
&\quad \mid \text{expr}\ \text{";"} \mid \text{printstmt}
\mid \text{";"}.
\end{aligned}
\]
Pemisahan ini memudahkan parser memilih cabang berdasarkan token
awal. Sebagai contoh, token \code{if} langsung memilih
\code{ifstmt}, sedangkan token \code{while} memilih
\code{whilestmt}.

\subsection{Ekspresi Berprioritas}

Ekspresi dipecah menjadi beberapa tingkat:
\[
\begin{aligned}
\text{expr} &\to \text{assign},\\
\text{assign} &\to \code{ID}\ \text{"="}\ \text{assign}
\mid \text{logor},\\
\text{logor} &\to \text{logand}\ \text{logor\_tail},\\
\text{logand} &\to \text{equality}\ \text{logand\_tail},\\
\text{equality} &\to \text{relation}\ \text{equality\_tail},\\
\text{relation} &\to \text{additive}\ \text{relation\_tail},\\
\text{additive} &\to \text{term}\ \text{additive\_tail},\\
\text{term} &\to \text{unary}\ \text{term\_tail},\\
\text{unary} &\to \text{"!"}\ \text{unary}
\mid \text{"-"}\ \text{unary}
\mid \text{"+"}\ \text{unary} \mid \text{primary}.
\end{aligned}
\]
Setiap tingkat menangani satu kelompok operator. Prioritas naik dari
\code{logor} ke \code{term}; unary dan primary berada di paling
dalam.

\subsection{Contoh Program}

Program \file{arith.sc} menguji prioritas dan tanda kurung:
\begin{lstlisting}[language=C]
int main() {
    int a = 2 + 3 * 4;
    int b = (2 + 3) * 4;
    print(a);
    print(b);
    return 0;
}
\end{lstlisting}
Ekspresi pertama menghasilkan AST dengan akar \code{+} dan anak
kanan \code{*}; ekspresi kedua menghasilkan akar \code{*} dengan
anak kiri \code{+}. Perbedaan ini berasal dari tanda kurung dan
tercermin dalam pohon urai.

\begin{warningbox}
Tata bahasa pada bab ini masih memuat rekursi kiri tidak langsung
pada beberapa aturan dan belum tentu LL(1). Bab IX akan menghapus
rekursi kiri dan melakukan left factoring sebelum parser top-down
dibangun.
\end{warningbox}
